\section{Step-by-step QEF vertex solve and payload}
\label{sec:qef}

\subsection{Why a QEF instead of an average}

A Surface Nets vertex is often placed at the average of the edge
crossings.  That is cheap but smooths sharp features: a box corner becomes
a rounded blob.  A QEF fits a point to the tangent planes of the crossings,
which preserves sharp edges better while still producing one vertex per
cell.  The engine regularizes the QEF toward the mass point so an
underconstrained cell does not explode.

\subsection{The twelve crossings}

For each edge $(i,j)$ with finite endpoints and opposite signs:
\begin{equation}
  t = \clip\left(\frac{d_i}{d_i-d_j},0,1\right),
  \quad
  \mathbf{p} = (1-t)\mathbf{p}_i + t\mathbf{p}_j,
  \quad
  \mathbf{n} = \operatorname{normalize}\bigl((1-t)\mathbf{n}_i+t\mathbf{n}_j\bigr).
\end{equation}
The corner normals $\mathbf{n}_i$ are the analytic gradients of the
trilinear field at the corner.  The edge normal is the interpolation of
the two corner normals, not the normalized edge direction.  This is what
makes the QEF preserve the smooth surface normal.

\subsection{Accumulation}

The normal equations are accumulated:
\begin{align}
  A &= \sum_e \mathbf{n}_e\mathbf{n}_e^{T} + \epsilon I,\\
  \mathbf{b} &= \sum_e \mathbf{n}_e(\mathbf{n}_e\cdot\mathbf{p}_e)
             + \epsilon\bar{\mathbf{p}}.
\end{align}
where $\bar{\mathbf{p}}$ is the average crossing and $\epsilon=10^{-2}$.
The matrix is symmetric positive semidefinite; the shift makes it
invertible.  Cramer's rule is used because it is branch-free and matches
the CPU's GLM evaluation order.  A determinant below $10^{-6}$ returns zero
and the caller falls back to the mass point.

\subsection{Corner snap}

If any corner SDF is exactly zero, the vertex is snapped to that corner:
\begin{equation}
  \mathbf{x} = C_i.
\end{equation}
Why?  Neighbouring cells sharing the corner would otherwise solve slightly
different QEF systems (their normals and crossings differ in the low bits)
and their vertices would drift apart, splitting the shared edge into two
one-sided edges.  Snapping to the exact corner makes every incident cell
produce the same bits.  This is a classic dual-contouring robustness fix
and it is essential for watertight terrain.

\subsection{Sentinel handling}

An edge with an infinite endpoint is skipped.  If no usable edge remains
but there are finite sign-changing edges blocked by a sentinel, the shader
averages the finite crossings instead of returning the cell centre.  The
fallback keeps the vertex on the real surface segment rather than floating
in the middle of the cell.

\subsection{Clamping}

The result is clamped to the cell:
\begin{equation}
  \mathbf{x} = \clip(\mathbf{x}, \mathbf{c}, \mathbf{c}+L).
\end{equation}
A QEF solution outside the cell would make the dual mesh non-manifold and
could place a vertex inside a neighbouring solid.  The clamp is the last
line of defence after regularization.

\subsection{Memoization}

The shader memoizes up to eight QEF positions per invocation.  A single
edge segment queries the same cells several times (owner selection, dedup,
rotation, fan).  The memo is a linear scan over eight entries; it is
value-identical to recomputing because the QEF is a pure function of the
node's eight corners and cube.  If the memo is full, the slot is chosen by
\code{id mod 8}; a collision simply evicts an older entry.

\subsection{The payload}

Stage 1 writes the final vertex from the node id:
\begin{lstlisting}[style=glsl,caption={Vertex payload layout.},label={lst:payload}]
uint o = (pc.vertexBase + slot) * 4u;
poolVerts[o+0] = vec4(p, 1.0);
poolVerts[o+1] = vec4(1.0, 1.0, data);
poolVerts[o+2] = vec4(n, uintBitsToFloat(uint(brush)));
poolVerts[o+3] = vec4(uintBitsToFloat(uint(type)), h);
\end{lstlisting}
The vertex is 64 bytes:
\begin{center}
\begin{tabular}{@{}c|l@{}}
\toprule
Word & Bytes\\
\midrule
0 & position.xyz (12) + color.x (4)\\
1 & color.yz (8) + texCoord.xy (8)\\
2 & normal.xyz (12) + brushIndex bits (4)\\
3 & textureType bits (4) + hsv (12)\\
\bottomrule
\end{tabular}
\end{center}
The \code{textureType} is stored as integer bits, not as a float value.
The vertex input attribute is \code{R32\_SINT}, and the ray-tracing path
uses \code{floatBitsToInt}; writing a numeric float would decode to a
garbage mapping mode and drop every non-triplanar vertex.

\subsection{One vertex per node}

Every surface node in the flattened table gets a vertex, including nodes
that the index stream never references.  This is intentional:
\begin{itemize}
\item it makes slot assignment a simple table order;
\item it makes the span size known from the host count;
\item it lets the finalize pass write one vertex per invocation without a
  compaction pass.
\end{itemize}
The wasted bytes are bounded by the number of surface nodes in the
neighbourhood and are small compared with the triangle stream.
